\section{Operational Lexicon}
\setstretch{1.15}
{\small

The following lexicon defines all terms, acronyms, and operational concepts used throughout this manuscript. Entries are organized thematically. Part VIII provides a complete acronym quick-reference index.

\subsection*{Part I: Core MDI Framework Terms}

\paragraph{Anti-Detention Covenant} The civil liberty guarantee codified in the Stewardship Authority charter. Guarantees unrestricted egress for all residents at all times following admission. Governs voluntary residency only. Distinct from existing California legal hold authority (5150 WIC, LPS Conservatorship), which is statutory law predating MDI.

\paragraph{ALMU (Asset Limited Modular Unit)} The precision-engineered 150 sq ft residential module within the MDI tower. Classified as a non-property interest. Includes independent wet bath, Murphy wall-bed, locking storage, induction cooktop, mini-refrigerator, and STC 65 acoustic door. Only the resident holds the key.

\paragraph{Automated-Human Operational Balance} The governing cybernetic principle. Automation applied to the mechanical layer (safety telemetry, robotic maintenance, logistics). Human judgment applied to the relational layer (Pod Steward welfare checks, peer specialist engagement, crisis response). Total automation and total human management are both rejected.

\paragraph{Bulk Majority Argument} The thesis that clearing 70 to 80 percent of the Pipeline C street population produces a population ecology shift, rendering the remaining voluntary nomadic Pipeline D population manageable via harm reduction rather than crisis management.

\paragraph{Dunbar Group (DG)} An ITP-assigned social cluster within a Dunbar Pod. Sizing is individualized: 5 to 8 individuals for socially sensitive residents; up to 20 to 30 for those with higher social tolerance. DG assignment is revisable as the resident stabilizes. Smaller DG is a clinical match, not a lower-status outcome.

\paragraph{Dunbar Pod} One of 13 independent residential partitions within the MDI tower. Approximately 154 units per pod, yielding 2,000 total. Four to five floors. Independent elevator loop. Time-slot-scheduled shared amenity access via RFID lockout. Derived from Dunbar's Number cognitive ceiling of approximately 150 for stable social relationships.

\paragraph{Efficiency Surplus} The net annual public financial benefit per tower generated through municipal service cost elimination. Best case: 82 million dollars per year. Most probable: 52.5 million dollars per year. Worst case: 33 million dollars per year. Capital recovery timeline: 29--71 months. Six-layer dividend: ER throughput, law enforcement reallocation, sanitation, commercial district recovery, community hospital relief, and front line worker conditions.

\paragraph{Exoskeletal Identity Structure} The functional role of a Pipeline C individual's accumulated hoard. Provides the physical boundary of a self-constructed world. Not random accumulation. Removing the hoard without replacing the identity structure triggers acute psychiatric destabilization. Intervention must be identity-preserving.

\paragraph{Tenancy Bridge Guarantee} The contractual mechanism ensuring the Stewardship Contract remains active as a legal guarantor of tenure until a new private market lease is verified, signed, and funded. No resident transitions to the independent market without a confirmed landing.

\paragraph{Coordinated Unit Activation} The simultaneous opening of all pre-matched tower units at One California Plaza. Transforms the housing offer from a generic referral into a specific, immediate, and irreversible commitment: named room, stored cart, kenneled pet, keycard issued at the moment of willingness.

\paragraph{Leaf Blower Effect} The consequence of encampment clearance without simultaneous housing placement. Produces geographic redistribution, not population reduction. ACT engagement relationships are severed. Engagement timelines reset to zero.

\paragraph{MDI (Material Dignity Infrastructure)} The end-to-end industrial pipeline for the engineered elimination of chronic street homelessness. Comprises Phase Zero (Metabolic Stabilization), Phase One (Field Architecture), Phase Two (Ground Floor Intake), Phase Three (Ontological Architecture), and five simultaneous operations.

\paragraph{Metabolic Stabilization (MS)} Phase Zero of the MDI pipeline. The 30 to 90 day clinical stage delivering sub-acute metabolic, neurological, and nutritional stabilization before any permanent housing placement occurs. The prerequisite layer that Housing First requires but the United States system does not currently provide.

\paragraph{MS Unit} The Metabolic Stabilization Unit. The clinical residential environment where Phase Zero stabilization occurs. Operated under MHRC licensure. Distinct from the ALMU residential track.

\paragraph{Nocturnal Access Protocol} Automated throttling of discretionary common floor access between 0000 and 0500 hours. Normalizes institutional social rhythm to match market-rate residential buildings. Medical, mental health, and safety access remain unrestricted globally.

\paragraph{Ontological Permanence Architecture (Ontological Security)} The physical and operational mechanisms providing structural predictability sufficient to eliminate the routine of being rousted (Giddens, 1991; Laing, 1960). Core markers: space for daily routines, sense of privacy, control over environment, coherent sense of self. The ALMU delivers all four simultaneously.

\paragraph{Pod Steward} A credentialed peer support specialist with lived experience of homelessness or psychiatric illness. One per Dunbar Pod of 154 individuals. Residential unit within the pod preferred. Primary function: recognition, not enforcement. Human welfare interpreter of life-safety telemetry. The relational layer of the Automated-Human Operational Balance.

\paragraph{Population Ecology Shift} The qualitative change in the street environment produced by substantially housing Pipeline C. Smaller remaining population. Less acute destabilization. Less visible public safety emergency. Diminished behavioral contagion. Changes Pipeline D management from crisis response to harm reduction.

\paragraph{Rapid Transition Protocol} The mechanism replacing eviction. Transfers a resident to a different pod or tower within 24 hours when behavior creates physical danger for pod members and de-escalation has failed. Medical necessity review precedes transfer. Resident retains Stewardship Contract. No resident becomes homeless again through behavioral crisis.

\paragraph{Sally Port (Pathway 1)} The formal, pre-matched intake pathway. Mobile crisis outreach and ACT team-facilitated warm offer to a named individual on the by-name HMIS registry. Unit, pod, and floor are pre-matched before the offer conversation begins. Highest-quality intake: the individual is known, pre-matched, and entering at a chosen moment of willingness.

\paragraph{Singular Prototype Threshold (SPT)} The eight binary verification metrics that must all pass at One California Plaza before any replication tower is acquired. A binary gate. Partial success does not authorize partial expansion. Metric 7 (Phase Zero clearance rate) is the causal test of the primary academic thesis.

\paragraph{Sovereign Acquisition via Receivership} The Stewardship Authority's mechanism for acquiring distressed commercial towers through court-supervised bankruptcy proceedings at verified floor valuations (typically 120 dollars per square foot; 70 to 85 percent below peak CRE valuations).

\paragraph{STC 65 Phase Zero Module} The physical module of each MHRC cluster. A 65-decibel sound enclosure designated as a state-owned clinical environment. Not a residential dwelling. Not a shelter. A therapeutic milieu unit analogous to a hospital step-down bed.

\paragraph{Stewardship Authority} The state-chartered quasi-governmental entity modeled on the Tennessee Valley Authority. Possesses delegated state police power and independent bonding capacity. Functions as MDI operator, CARE Court petitioner, real estate acquirer, and outcome verifier. Addresses the coordination failure that paralyzes existing Los Angeles homelessness governance.

\paragraph{Stewardship Contract} The legal instrument governing residential tenure within the tower. Classified as a non-property interest. Resident holds occupancy, privacy, and non-interference rights equivalent to a residential lease without triggering housing court eviction jurisdiction.

\paragraph{Hygiene Intake Kit} The gift offered to Diogenes-pattern individuals at hygiene intake. Framed as addition, not replacement. Staff bag and tag soiled clothing pending individual consent for disposal.

\paragraph{Resident Civic Engagement Pathway} Voluntary participation tracks transitioning residents from passive service recipients to active stakeholders. Includes Garden Stewardship, Robotic Swarm Supervision, Community Kitchen Coordination, and Digital Access Center Support. Earns Stewardship Status, a social recognition category within the pod.

\paragraph{Three Ps Infrastructure} The structural elimination of three primary shelter refusal barriers: Pets (on-site facility, veterinary node, dog run); Partners (co-location within same Dunbar Pod); Possessions (photographed, GPS-tagged, RFID-accessible secure storage).

\paragraph{Warm Offer} A specific, immediate, irreversible housing offer made at the moment of the individual's first window of willingness following an acute crisis event. Pre-confirmed unit, kenneled pet, stored possessions, completed intake documentation. Eliminates bureaucratic delay and every rational refusal barrier simultaneously.

\subsection*{Part II: Population Pipelines}

\paragraph{Pathway 3 (Secondary Clinical Transfer)} Individual transferred following field contact, non-criminal law enforcement determination, or ER discharge to a lower-acuity setting. Uses the rear-wall Zone A entrance, physically separated from the open commons. Prevents crisis contagion in the resource center population.

\paragraph{Pipeline A} The near-homeless and voluntarily transitioning. No major psychiatric comorbidity. Engages voluntarily with the Coordinated Entry System. Rapid placement track. CES is the existing instrument; MDI accelerates it by creating surplus inventory.

\paragraph{Pipeline B} The encamped but engageable. Long-term encampment with intact insight. Refusal directed at congregate shelter barriers (pets, partners, possessions), not housing itself. Warm offer at acute crisis window. 58 to 70 percent acceptance rate with a specific, immediate offer.

\paragraph{Pipeline C} The behaviorally calcified chronic. Long Duration of Untreated Psychosis (5 to 15 years). Active schizophrenia spectrum disorder. Requires ACT sustained presence 12 to 24 months plus the legal lever system. Primary driver of public safety concern and street visibility.

\paragraph{Pipeline C Riparian Sub-Variant} The hidden riparian population. Mobile, terrain-adaptive, encamped in environmentally sensitive sites (LA River, Ballona Creek, storm drains). Estimated 20 to 30 percent above PIT count Pipeline C figure due to systematic undercount. Requires environmental enforcement track (field mapping + CWA abatement) paired with simultaneous warm offer.

\paragraph{Pipeline D} The voluntary nomadic. Not metabolically destabilized in the Pipeline C clinical sense. Organized life around outdoor community, informal economy, and deliberate rejection of institutional structures. Instrument: harm reduction without housing requirement. No legal instrument applies. Clearance of Pipeline C changes the street environment sufficiently that Pipeline D becomes manageable via harm reduction alone.
\subsection*{Part III: Clinical Terms}

\paragraph{Anosognosia} A neurological incapacity to perceive one's own psychiatric condition. Up to 50 percent of schizophrenia-spectrum individuals present with it. The frontal lobe damage producing psychosis simultaneously destroys the self-monitoring faculty. Outreach strategies premised on persuasion target a faculty already destroyed by the disease.

\paragraph{Behavioral Calcification} The progressive rigidity of survival routines and threat-perception systems produced by years of untreated psychosis and street-level chronic stress. Key barrier distinguishing Pipeline C from Pipeline B. Duration directly correlates with intervention resistance.

\paragraph{Day-30 Functional Indicators} The five ITP measurable goals used as the Phase Zero clearance threshold: (1) sustained hygiene maintenance without prompting; (2) six or more hours sleep without pharmacological support; (3) orientation to person, place, and date across three consecutive days; (4) one structured group activity per day; (5) one independent ADL task per day. Progress on two or more equals clearance.

\paragraph{Diogenes Syndrome} Extreme self-neglect, domestic squalor, social withdrawal, and adamant refusal of help. The accumulated hoard functions as an Exoskeletal Identity Structure. Intervention must be identity-preserving. Hoard removal without identity structure replacement triggers acute destabilization.

\paragraph{DSM-5 (Diagnostic and Statistical Manual of Mental Disorders, Fifth Edition)} The APA diagnostic classification system governing all provisional and finalized diagnoses throughout the Phase Zero Three-Phase Diagnostic Protocol.

\paragraph{DUP (Duration of Untreated Psychosis)} The accumulation of years of untreated psychosis producing measurable progressive brain changes: increasing cognitive rigidity, treatment resistance, and functional decline. Pipeline C typically presents 5 to 15 years DUP. Standard clinical timescales do not apply to this cohort.

\paragraph{Hyponatremia} Potentially fatal cerebral edema produced by rapid water loading in severely dehydrated individuals with depleted sodium. Life-safety constraint on the rehydration protocol: WHO-formulated electrolyte solution is required before water. Governs the robotic kitchen and pharmacy protocols.

\paragraph{Refeeding Syndrome} Life-threatening electrolyte shifts (cardiac arrhythmia, neurological complications) produced by rapid carbohydrate reintroduction in severely malnourished individuals. Clinical triage layer required at intake. BMI below 16, visible severe wasting, or reported days without food triggers the Gradual Reintroduction Protocol before open cafeteria access.

\paragraph{Three-Phase Diagnostic Protocol} Phase 1 (Hours 0 to 72): Provisional DSM-5 diagnosis by LPHA; billable as Mental Health Services Assessment. Phase 2 (Days 1 to 14): Collateral gathering, daily progress notes, sustained billing trail. Phase 3 (By Day 30): Finalized DSM-5 diagnosis signed by LPHA; ITP signed; full Medi-Cal billing authorized.

\paragraph{Wernicke Encephalopathy} Neurological emergency produced by Thiamine (Vitamin B1) deficiency during nutritional refeeding. Clinical contraindication to rapid carbohydrate reintroduction in severely malnourished arrivals. Governed by the Refeeding Syndrome clinical triage layer.

\subsection*{Part IV: Legal and Regulatory Terms}

\paragraph{5150 Hold} California Welfare and Institutions Code Section 5150. Authorizes a 72-hour involuntary psychiatric hold for individuals posing danger to self, danger to others, or meeting grave disability criteria. Predates MDI. Provides the acute triage authority during the crisis window. Not an MDI-created authority.

\paragraph{AOT (Assisted Outpatient Treatment)} Court-ordered outpatient treatment for individuals with documented history of treatment non-compliance and repeated psychiatric hospitalizations, incarcerations, or threats/acts of violence. Governed by Laura's Law (2002). Mandates medication compliance and outpatient engagement. Non-compliance triggers court review and potential LPS escalation.

\paragraph{ARO (Adaptive Reuse Ordinance)} The Los Angeles ordinance (February 2026 expansion) eliminating minimum unit size requirements and public hearing mandates for qualifying commercial-to-residential conversions. The MDI ALMU at 150 sq ft qualifies under the expanded ARO.

\paragraph{CARE Court} Community Assistance, Recovery, and Empowerment Act (2022). California court mechanism for adults with schizophrenia spectrum or other psychotic disorders. Non-custodial. Establishes a 12 to 24 month court-monitored CARE Agreement including housing, medication, and wraparound services. The Stewardship Authority petitions as a qualifying entity using by-name HMIS registry documentation.

\paragraph{CWA (Clean Water Act, 1972)} Federal statute providing the legal basis for environmental abatement of riparian encampments. Fecal coliform contamination exceeding NPDES permit standards and biological waste infiltration into protected habitat constitute documented violations providing abatement authority independent of housing law.

\paragraph{Federal Supremacy Leasing} Legal mechanism invoking federal preemption of conflicting municipal regulations for tower space occupied by federal agencies or federally funded programs. Creates a secondary legal shield for federal service co-location floors against municipal obstruction.

\paragraph{Grants Pass Framework} \textit{City of Grants Pass v. Johnson} (2024). The Supreme Court ruling governing municipal authority over encampment clearance. Riparian abatement is explicitly outside this framework, as the legal basis is environmental protection (CWA, Army Corps jurisdiction), not anti-camping ordinances.

\paragraph{Home Rule Stabilization Ordinance} State ordinance pre-empting conflicting local land use regulations for qualifying projects. Codifies state housing law supremacy into automatic by-right conversion authority for commercial structures acquired through the Sovereign Acquisition process.

\paragraph{IMD Exclusion (Institution for Mental Disease)} The federal Medicaid rule excluding federal financial participation for mental health facilities with more than 16 beds. Resolved architecturally by the sub-16-unit MHRC cluster structure: each cluster falls below the threshold and files independently. No waiver required.

\paragraph{LPS Conservatorship (Lanterman-Petris-Short Act)} Highest-restriction legal instrument. Court appoints Public Guardian or designated conservator for individuals deemed gravely disabled and unable to provide for food, clothing, or shelter due to mental health disorder or severe substance use disorder (SB 43 expansion, 2023). Conservator may consent to housing placement and medical treatment. Subject to periodic court review.

\paragraph{Medical Classification Waiver} DHCS petition enabling ALMU residential units to be classified as community-based residential treatment for Medi-Cal billing purposes. Preserves residential character while enabling clinical billing that funds the field architecture.

\paragraph{NPDES (National Pollutant Discharge Elimination System)} Clean Water Act permitting system. Fecal coliform contamination of LA waterways from encampments exceeds NPDES permit standards, constituting documented federal regulatory violations and the legal basis for environmental abatement proceedings.

\paragraph{SB 330} State law vesting zoning entitlements at the date of application for qualifying projects. Prevents subsequent zoning amendments from altering the entitlement after the Stewardship Authority files a conversion application.

\paragraph{Stewardship Agreement} The resident's contractual instrument under the Stewardship Contract. Grants occupancy, privacy, and non-interference rights. Classified as a non-property interest to bypass housing court jurisdiction while preserving residential dignity protections.

\paragraph{TVA Model (Tennessee Valley Authority)} The template for the Stewardship Authority. Quasi-governmental, state-chartered, operationally autonomous from municipal government. Consolidates field operations, legal petitioning, real estate acquisition, management, and outcome verification under a single mandate-limited entity.

\paragraph{WIC (Welfare and Institutions Code)} California statute. Section 5150 provides the 72-hour involuntary psychiatric hold authority applied in the Three-Gate Triage Decision Architecture as the governing authority for the Emergency Room pathway.

\subsection*{Part V: Operational and Field Architecture Terms}

\paragraph{ACT (Assertive Community Treatment)} Evidence-based intensive outreach model. Team composition: mobile psychiatrist (shared), RN, licensed social worker, outreach case manager, peer support specialist, substance use counselor, team coordinator (7-8 FTE). Caseload ratio: 80-120 Pipeline C clients. Medi-Cal-billable in California at $150 to $190 per day per enrolled client.

\paragraph{By-Name HMIS Registry} The operational backbone of the Field Architecture. Every individual in the target zone registered by name in an HMIS instance managed by the Stewardship Authority. Every housing unit pre-matched to a named individual before tower opening. The warm offer is specific because the registry preparation was specific.

\paragraph{CES (Coordinated Entry System)} The existing Los Angeles coordinated intake and assessment system for homeless services. Pipeline A individuals engage voluntarily through CES. MDI accelerates CES by creating surplus inventory.

\paragraph{CIT (Crisis Intervention Team)} Specialized law enforcement co-responder model. 40-hour specialized training. Mental health clinician embedded with CIT officers. Routes individuals directly into designated psychiatric emergency facilities, bypassing criminal jail. Validates dual-role crisis node staffing in Zone A and Zone B.

\paragraph{CPTED (Crime Prevention Through Environmental Design)} Architectural alternative to camera-dense digital surveillance. Designs spaces enabling residents to naturally observe common areas as part of daily movement. Glass-partitioned common rooms, high-visibility stairwells, digital access center, and pod kitchen at circulation chokepoints.

\paragraph{EmPATH (Emergency Psychiatric Assessment Treatment and Healing)} Non-coercive psychiatric emergency unit design model. Open floor plan, recliner chairs, residential aesthetic, rapid stabilization, de-escalation as environmental design first response. Validates Zone A and Zone B ground floor design philosophy.



\paragraph{Outreach Workers} MDI field assessment personnel. Jointly with ACT teams and law enforcement, apply the Three-Gate Triage Decision Architecture to every crisis presentation.

\paragraph{HMIS (Homeless Management Information System)} HUD-required data system for tracking homeless individuals and services. The MDI Stewardship Authority operates a proprietary HMIS instance as the by-name registry backbone.

\paragraph{LAHSA (Los Angeles Homeless Services Authority)} The joint authority of the City and County of Los Angeles responsible for coordinating and implementing homeless services. Administers the annual Point-in-Time count. Source data for MDI field architecture sizing.

\paragraph{LARWQCB (Los Angeles Regional Water Quality Control Board)} State regulatory body governing water quality in the Los Angeles region. Enforcement records document Clean Water Act violations produced by riparian encampments. Provides the legal basis for environmental abatement independent of housing law.

\paragraph{PIT Count (Point-in-Time Count)} The annual one-night survey of the homeless population mandated by HUD. Systematically undercounts the most hidden encampments, particularly Pipeline C Riparian individuals. 2025 Greater Los Angeles PIT: approximately 75,312 total homeless, 58 percent unsheltered.

\subsection*{Part VI: Architectural and Physical Design Terms}

\paragraph{Residential Data Minimization Protocol} The only data exported from a residential ALMU to the central cybernetic grid is a binary safe/not-safe occupancy indicator. All other environmental data processed locally and purged within 24 hours. Structurally precludes behavioral profiling and the Panopticon Effect.

\paragraph{Biophilic Infrastructure} Two-tier system. Tier 1: 200-300 sq ft planted pod-level alcove on every 4-5 floors, adjacent to pod common kitchen. Tier 2: Tower-level destination sanctuary every 7 floors with hydroponic garden, oxygenated lounge, rooftop greenspace. Reduces cortisol, improves mood regulation, decreases agitation incidents.

\paragraph{Real-time Unit Availability Display} Live display of unit and resource availability across the network. Deployed in the ground floor Resource Center and Transportation Asset Management system. Neutralizes scarcity-driven anxiety before it becomes behavioral agitation.

\paragraph{Digital Access Center} Low-distraction, high-bandwidth vocational and educational access hub. Gigabit internet, hardware lending, vocational training, legal self-defense resources (benefits enrollment, ID restoration, credit rehabilitation), gig-work platform access.

\paragraph{HEPA (High Efficiency Particulate Air)} Air filtration standard. Clinical-grade HEPA scrubbing deployed in ALMU units and acuity-segmented hygiene modules. Air quality exceeds average Los Angeles apartment standards.

\paragraph{HLM (Habitation \`{a} Loyer Mod\'{e}r\'{e})} French social housing designation. Large-scale high-rise social housing complexes where the Gardien/Pod Steward model was developed and validated.

\paragraph{HVAC (Heating, Ventilation, and Air Conditioning)} Individual HVAC control within a 68 to 78 degree operating band in each ALMU. Tower-level Tier 2 biophilic sanctuary operates at 50 percent of original HVAC design load surplus, producing elevated oxygen concentration.

\paragraph{MAP (Managed Alcohol Program)} Medically supervised alcohol management facility for individuals in active, functionally incapacitating heavy substance dependence. Pre-ALMU track for this sub-population. Post-MAP stabilization, transfer to the ALMU track follows.

\paragraph{MHRC (Mental Health Rehabilitation Center)} A 24-hour DHCS-licensed facility classification for adults with mental disorders providing intensive rehabilitative services for sub-acute stabilization, skill-building, and independence. The MDI Phase Zero cluster model. Sub-16-unit threshold resolves the IMD Exclusion architecturally.

\paragraph{MTBF (Mean Time Between Failures)} Equipment reliability standard. Pressurized sewage ejectors at One California Plaza carry a 50,000-hour MTBF mandate.

\paragraph{RFID (Radio Frequency Identification)} Electronic identification token system. Used for pod boundary enforcement on shared amenity floors via time-slot scheduling, secure possession storage access, and bicycle transportation asset tracking.

\paragraph{STC (Sound Transmission Class)} Measure of sound attenuation performance. STC 55 is the biological floor (prevents hearing damage). STC 65 is the MDI ontological threshold, which reduces ambient sound transfer to clinical quiet-room standards, neutralizing hypervigilance triggers in trauma-impacted residents.

\paragraph{UV-C (Ultraviolet-C)} Germicidal ultraviolet light wavelength. UV-C scrubbing deployed in acuity-segmented hygiene modules and ALMU HVAC systems for clinical-grade disinfection.

\paragraph{VOC (Volatile Organic Compound)} Airborne chemical compounds released by building materials. MDI biophilic node species selected for low-VOC properties, drought tolerance, and sensory appropriateness.

\subsection*{Part VII: Fiscal Terms}

\paragraph{CRE (Commercial Real Estate)} The distressed commercial real estate providing the MDI acquisition target. One California Plaza: 120 dollars per square foot floor valuation (approximate 80 percent below peak CRE valuations). The 16 million square foot network conversion opportunity across the distressed downtown corridor.

\paragraph{DHCS (Department of Health Care Services)} California state agency. Licenses MHRC facilities via Form 1813. Governs Medi-Cal specialty mental health billing standards. MDI petitions DHCS for the Medical Classification Waiver.

\paragraph{FMAP (Federal Medical Assistance Percentage)} The federal cost-sharing rate for Medi-Cal expenditures. Approximately 50 percent federal financial participation on MHRC sub-acute specialty mental health billing. The load-bearing fiscal base of Phase Zero.

\paragraph{Four-Layer Funding Stack} Layer 0: MHRC Medi-Cal sub-acute billing (Phase Zero). Layer 1: ACT Medi-Cal billing (clinical services at 150 to 190 dollars per day per enrolled client). Layer 2: Measure Alpha and MHSA Full Service Partnership (field architecture, Pod Stewards, residential management). Layer 3: Efficiency Surplus (capital recovery through municipal service cost elimination). No layer depends on speculative administrative approval for core viability.

\paragraph{Measure Alpha} Los Angeles County measure generating 843 million dollars annually for homelessness services. The MDI Stewardship Authority operationalizes this funding by providing a verifiable infrastructure destination with measurable per-resident outcome verification.

\paragraph{MHSA (Mental Health Services Act)} California state act generating approximately 3.5 billion dollars statewide annually. Los Angeles County allocation: approximately 400 million dollars per year. ACT teams for the chronically homeless unsheltered are among the highest priority eligible expenditures under the Full Service Partnership category.

\paragraph{SAMHSA (Substance Abuse and Mental Health Services Administration)} Federal agency funding ACT implementation through competitive grants of \$500K--\$2M per program.

\bigskip
\subsection*{Part VIII: Acronym Quick-Reference Index}

{\singlespacing\small
\begin{tabularx}{\textwidth}{@{}p{1.5cm}Xp{1.9cm}X@{}}
\toprule
\rowcolor{rowgray}
\textbf{Acronym} & \textbf{Expansion} & \textbf{Acronym} & \textbf{Expansion} \\
\midrule
ACT      & Assertive Community Treatment              & ITP     & Individualized Treatment Plan \\
\rowcolor{rowgray}
ALMU     & Asset Limited Modular Unit                 & LAHSA   & Los Angeles Homeless Services Authority \\
AOT      & Assisted Outpatient Treatment              & LARWQCB & LA Regional Water Quality Control Board \\
\rowcolor{rowgray}
ARO      & Adaptive Reuse Ordinance                   & LPS     & Lanterman-Petris-Short Act \\
BMI      & Body Mass Index                            & LPHA    & Licensed Practitioner of the Healing Arts \\
\rowcolor{rowgray}
CARE Court & Community Assistance, Recovery, and Empowerment Act Court & MAP & Managed Alcohol Program \\
CES      & Coordinated Entry System                   & MDI     & Material Dignity Infrastructure \\
\rowcolor{rowgray}
CIT      & Crisis Intervention Team                   & MHRC    & Mental Health Rehabilitation Center \\
CPTED    & Crime Prevention Through Environmental Design & MHSA & Mental Health Services Act \\
\rowcolor{rowgray}
CRE      & Commercial Real Estate                     & MS      & Metabolic Stabilization \\
CWA      & Clean Water Act                            & MTBF    & Mean Time Between Failures \\
\rowcolor{rowgray}
DG       & Dunbar Group                               & NIMBY   & Not In My Backyard \\
DHCS     & Department of Health Care Services         & NPDES   & National Pollutant Discharge Elimination System \\
\rowcolor{rowgray}
DSM-5    & Diagnostic and Statistical Manual, 5th Ed. & PIT    & Point-in-Time (Count) \\
DUP      & Duration of Untreated Psychosis            & REM     & Rapid Eye Movement \\
\rowcolor{rowgray}
EmPATH   & Emergency Psychiatric Assessment Treatment and Healing & RFID & Radio Frequency Identification \\
FMAP     & Federal Medical Assistance Percentage      & SAMHSA  & Substance Abuse and Mental Health Services Administration \\
\rowcolor{rowgray}
         &                                            & SPT     & Singular Prototype Threshold \\
FTE      & Full-Time Equivalent                       & STC 65  & Sound Transmission Class 65 Module \\
\rowcolor{rowgray}
HEPA     & High Efficiency Particulate Air            & STC     & Sound Transmission Class \\
HLM      & Habitation \`{a} Loyer Mod\'{e}r\'{e}     & TVA     & Tennessee Valley Authority \\
\rowcolor{rowgray}
HMIS     & Homeless Management Information System     & UV-C    & Ultraviolet-C \\
HVAC     & Heating, Ventilation, and Air Conditioning & VOC     & Volatile Organic Compound \\
\rowcolor{rowgray}
IMD      & Institution for Mental Disease             & WIC     & Welfare and Institutions Code \\
\bottomrule
\end{tabularx}
}

